% Generated draft: requires mathematical and source review.
\documentclass[11pt]{article}
\usepackage[margin=25mm]{geometry}
\usepackage{fontspec}
\setmainfont{Latin Modern Roman}
\usepackage{amsmath,amssymb,mathtools,unicode-math}
\setmathfont{Latin Modern Math}
\usepackage{enumitem,needspace,xurl,hyperref}
\setlength{\parindent}{0pt}
\setlength{\parskip}{5pt}
\setlength{\emergencystretch}{3em}
\newcommand{\N}{\mathbb{N}}
\newcommand{\Z}{\mathbb{Z}}
\newcommand{\Q}{\mathbb{Q}}
\newcommand{\R}{\mathbb{R}}
\newcommand{\C}{\mathbb{C}}
\newcommand{\F}{\mathbb{F}}
\newcommand{\A}{\mathbb{A}}
\newcommand{\PP}{\mathbb{P}}
\newcommand{\E}{\mathbb{E}}
\newcommand{\eps}{\varepsilon}
\newcommand{\dd}{\,\mathrm{d}}
\begin{document}
\section*{UnsolvedMath ID: 30003451}
\textbf{Source problem code:} catalogue entry
\subsection{Definitions and mathematical statement}
The supplied record contains a conjecture and a separate question about evaluation vectors. It does not define the notation sufficiently to identify a unique mathematical formulation.

For clarity, the following definitions describe a conditional homogeneous-form interpretation. Let \(n\geq1\) and \(d\geq0\) be integers, and let \(V\) be a finite-dimensional real vector space of homogeneous polynomials of degree \(2d\) in \(n\) variables. For a basis \(b_1,\ldots,b_m\) of \(V\) and a nonzero vector \(v\in\R^n\), define
\[
 s_V(v)=(b_1(v),\ldots,b_m(v))\in\R^m.
\]
For \(f\in V\), its real projective zero set is
\[
 Z_{\R\PP^{n-1}}(f)=\{[v]\in\R\PP^{n-1}:f(v)=0\}.
\]
Homogeneity makes this zero set well defined. Replacing a representative \(v\) by \(tv\), where \(t\neq0\), multiplies its evaluation vector by \(t^{2d}\); hence linear independence of evaluation vectors is independent of the chosen representatives and of the basis of \(V\).

For such a space, one possible moment cone is
\[
 \mathcal M(V)=\left\{\sum_{j=1}^{r}w_j s_V(v_j): r\in\Z_{\geq0},\ w_j>0,\ v_j\in\R^n\setminus\{0\}\right\}.
\]
Its Carath\'{e}odory number is the least nonnegative integer \(C\), if one exists, such that every element of \(\mathcal M(V)\) has such a representation with at most \(C\) summands. A corresponding maximal finite projective-zero count would be
\[
 \max\left\{\#Z_{\R\PP^{n-1}}(f): f\in V,\ f(x)\geq0\ \text{for all }x\in\R^n,\ \#Z_{\R\PP^{n-1}}(f)<\infty\right\},
\]
provided this maximum exists. These definitions explain a possible interpretation only: the supplied record does not identify the space, cone, or class of forms used to define its \(C_{2d}\) and \(\beta(2d)\), or confirm that its notation \(B_{n,2d}\) denotes this kind of space.

The conjecture as supplied is that, for every integer \(d\geq3\),
\[
 \beta(2d)\leq C_{2d}\leq\beta(2d)+1.
\]
Its numerical meaning remains unresolved until the definitions of both quantities and their dependence on the number of variables are specified.

The accompanying question asks whether, whenever \(f\in B_{n,2d}\) is nonnegative on \(\R^n\) and has finite projective zero set consisting of distinct points
\[
 \{z_1,\ldots,z_k\},
\]
the vectors
\[
 s_{B_{n,2d}}(z_1),\ldots,s_{B_{n,2d}}(z_k)
\]
are linearly independent over \(\R\). Under the conditional interpretation \(B_{n,2d}=V\) above, choose nonzero representatives \(v_i\) of \(z_i\). The question is then whether
\[
 \left(\sum_{i=1}^{k}a_i g(v_i)=0\ \text{for every }g\in V\right)
 \quad\Longrightarrow\quad a_1=\cdots=a_k=0
\]
for every \((a_1,\ldots,a_k)\in\R^k\). If \(k=0\), the empty family is linearly independent by convention. The source does not specify the ranges of \(n,d\) for this separate question; its explicit restriction \(d\geq3\) belongs to the conjecture. The projective space being real and the interpretation of \(B_{n,2d}\) as a homogeneous polynomial space require confirmation from the source.
\subsection{Short English statement}
The supplied conjecture places the Carath\'{e}odory number between a maximal finite projective-zero count and that count plus one, for \(d\geq3\). A separate question asks whether the evaluations at all finitely many projective zeros of a nonnegative form are linearly independent. The defining spaces, parameter ranges, and meanings of several symbols require source confirmation. The supplied open status is dated curation, not an independently verified conclusion.
\subsection{Sources}
\begin{itemize}
\item Local registry: the supplied complete source record for permanent UnsolvedMath ID \(30003451\), titled ``Carath\'{e}odory Bounds and Independent Evaluation Vectors,'' in the Oberwolfach Reports Open Problems collection. Exact source problem code: \texttt{catalogue entry}. The record gives proposed year 2017 and names no individual proposer.
\item Supplied report citation: \emph{Real Algebraic Geometry With a View Toward Moment Problems and Optimization} (2018), Oberwolfach Reports, pp.\ 771--862, DOI \texttt{10.4171/owr/2017/14}. Original Oberwolfach report: \url{https://doi.org/10.4171/owr/2017/14}. The supplied background describes these reports as documenting mathematical workshops and their open-problem discussions.
\item Supplied literature reference: \emph{Positive Gorenstein ideals}, \url{https://arxiv.org/abs/1203.3775}. The record supplies no author or precise passage connecting this reference to the questions.
\item Literature assessment dated exactly \texttt{2026-08-21}: the supplied curation labels the problem open and reports that the report establishes upper bounds and examples relating the Carath\'{e}odory number to the maximal finite projective-zero count, while the displayed conjecture and accompanying general question remain unresolved. This is supplied dated curation derived from the OWR record, not independently verified here; its status and cited claims require independent checking.
\item The registry labels the statement \texttt{corrected\_verified} and says its mathematical content was preserved while notation, typography, and sentence structure were repaired for publication. This is a supplied editorial qualification and does not resolve the missing definitions or constitute mathematical review of this draft.
\end{itemize}
\end{document}
